% !TEX root = ../main.tex
%% Corrigé du devoir surveillé n° 1 — vérifié par verifications/eval_ds.py.

\corrigeSujetEval{DS1}{Devoir surveillé \numero1}

\begin{corrigeEval}{5 points}{Ensembles et applications}
\begin{questions}
  \item On met $f$ sous forme canonique : $f(x)=(x-1)^{2}-1$. Pour
    $x\in\intFF{-1}{2}$, $x-1$ décrit $\intFF{-2}{1}$, donc $(x-1)^{2}$ décrit
    $\intFF{0}{4}$ et $\res{f\bigl(\intFF{-1}{2}\bigr)=\intFF{-1}{3}}$. Ensuite
    \[0\leqslant f(x)\leqslant3\iff1\leqslant(x-1)^{2}\leqslant4
      \iff1\leqslant|x-1|\leqslant2,\]
    d'où $\res{f^{-1}\bigl(\intFF{0}{3}\bigr)=\intFF{-1}{0}\cup\intFF{2}{3}}$.
    \pts{1,5 pt}
  \item $f(0)=f(2)=0$ : $f$ est \res{\text{non injective}} ; $f\geqslant-1$,
    donc $-2$ n'a pas d'antécédent : $f$ est \res{\text{non surjective}}. Sur
    $\intFO{1}{+\infty}$, $f$ est continue, strictement croissante, avec
    $f(1)=-1$ et $f\to+\infty$ : elle induit une bijection $g$ sur
    $\intFO{-1}{+\infty}$. Pour $y\geqslant-1$, l'équation $(x-1)^{2}=1+y$
    d'inconnue $x\geqslant1$ donne $\res{g^{-1}(y)=1+\sqrt{1+y}}$.
    \pts{1,5 pt}
  \item Si $y\in h(A\cap B)$, alors $y=h(x)$ avec $x\in A$ et $x\in B$, donc
    $y\in h(A)\cap h(B)$. Réciproquement, si $h$ est injective et
    $y\in h(A)\cap h(B)$, on écrit $y=h(a)=h(b)$ avec $a\in A$, $b\in B$ ;
    l'injectivité donne $a=b\in A\cap B$, donc $y\in h(A\cap B)$. Avec $f$,
    $A=\ens{0}$ et $B=\ens{2}$ :
    $\res{f(A\cap B)=\varnothing\subsetneq\ens{0}=f(A)\cap f(B)}$.
    \pts{2 pts}
\end{questions}
\begin{center}\begin{minipage}{\linewidth}\centering
\begin{tikzpicture}[x=0.95cm,y=0.62cm,line cap=round]
  \fill[otGreenDark!10] (-1.6,0) rectangle (3.6,3);
  \draw[otGreenDark!60,line width=0.4pt,dash pattern=on 2pt off 1.5pt] (-1.6,3) -- (3.6,3);
  \draw[otMuted,line width=0.5pt,-{Stealth[length=1.5mm]}] (-1.6,0) -- (3.8,0);
  \draw[otMuted,line width=0.5pt,-{Stealth[length=1.5mm]}] (0,-1.5) -- (0,4.2);
  \draw[otNavy,line width=1.1pt] plot[domain=-1.35:3.35,samples=80] ({\x},{\x*\x-2*\x});
  \draw[otOrange,line width=2.4pt] (-1,0) -- (0,0);
  \draw[otOrange,line width=2.4pt] (2,0) -- (3,0);
  \foreach \xx in {-1,1,2,3}
    \node[font=\scriptsize,anchor=north] at (\xx,-0.08) {$\xx$};
  \node[font=\scriptsize,anchor=north east] at (-0.02,-0.08) {$0$};
  \node[font=\scriptsize,anchor=east] at (-0.05,3) {$3$};
  \node[font=\scriptsize,anchor=east] at (-0.05,-1) {$-1$};
  \draw[otMuted,line width=0.4pt,dash pattern=on 1.5pt off 1.5pt] (1,0) -- (1,-1) -- (0,-1);
  \node[font=\scriptsize,text=otNavy,anchor=west] at (3.3,3.9) {$y=x^{2}-2x$};
\end{tikzpicture}\\*[0.5mm]
\legende{L'image réciproque de la bande $0\leqslant y\leqslant3$ (en vert) :
la réunion des deux segments orange $\intFF{-1}{0}\cup\intFF{2}{3}$.}
\end{minipage}\end{center}
\end{corrigeEval}

\begin{corrigeEval}{5 points}{Sommes et récurrences}
\begin{questions}
  \item $k\cdot k!=(k+1)\,k!-k!=(k+1)!-k!$ : la somme est télescopique,
    \[\sum_{k=1}^{n}k\cdot k!=\sum_{k=1}^{n}\bigl((k+1)!-k!\bigr)=\res{(n+1)!-1}.\]
    \pts{1 pt}
  \item On fait apparaître $(n+1)!$ :
    \[\frac{1}{k+1}\binom nk=\frac{n!}{(k+1)!\,(n-k)!}
      =\frac{1}{n+1}\cdot\frac{(n+1)!}{(k+1)!\,(n-k)!}=\frac{1}{n+1}\binom{n+1}{k+1}.\]
    Avec $j=k+1$ et la formule du binôme (il manque le terme $j=0$) :
    \[\sum_{k=0}^{n}\frac{1}{k+1}\binom nk=\frac{1}{n+1}\sum_{j=1}^{n+1}\binom{n+1}{j}
      =\res{\frac{2^{n+1}-1}{n+1}}.\]
    \pts{1,5 pt}
  \item On somme d'abord sur $i$, à $j$ fixé :
    \[\sum_{1\leqslant i\leqslant j\leqslant n}i=\sum_{j=1}^{n}\frac{j(j+1)}{2}
      =\frac12\Bigl(\frac{n(n+1)(2n+1)}{6}+\frac{n(n+1)}{2}\Bigr)
      =\res{\frac{n(n+1)(n+2)}{6}}.\]
    \pts{1,5 pt}
  \item Pour $n=1$, c'est une égalité. Si $(1+x)^{n}\geqslant1+nx$, on
    multiplie par $1+x\geqslant0$ (c'est là que sert $x\geqslant-1$) :
    \[(1+x)^{n+1}\geqslant(1+nx)(1+x)=1+(n+1)x+nx^{2}\geqslant1+(n+1)x.\]
    \pts{1 pt}
\end{questions}
\end{corrigeEval}

\begin{corrigeEval}{4 points}{Trigonométrie}
\begin{questions}
  \item Avec $s=\sin x$ et $\cos^{2}x=1-s^{2}$ :
    \[2(1-s^{2})+3s-3=0\iff2s^{2}-3s+1=0\iff s\in\Bigl\{1,\frac12\Bigr\}.\]
    Sur $\intFO{0}{2\pi}$ : $\res{x\in\Bigl\{\dfrac\pi6,\dfrac\pi2,\dfrac{5\pi}6\Bigr\}}$.
    \pts{1,5 pt}
  \item Formule de duplication et signe : $\cos^{2}\dfrac\pi8=\dfrac{1+\cos\frac\pi4}{2}=\dfrac{2+\sqrt2}{4}$
    et $\cos\dfrac\pi8>0$, donc $\res{\cos\dfrac\pi8=\dfrac{\sqrt{2+\sqrt2}}{2}}$.
    \pts{1 pt}
  \item Par la formule de Moivre :
    \[\cos3\theta=\operatorname{Re}\bigl((\cos\theta+\mathrm{i}\sin\theta)^{3}\bigr)
      =\cos^{3}\theta-3\cos\theta\sin^{2}\theta=4\cos^{3}\theta-3\cos\theta.\]
    Avec $t=\cos\theta$, $8t^{3}-6t-1=2\cos3\theta-1$ s'annule dès que
    $\cos3\theta=\frac12$, en particulier pour
    $\theta\in\bigl\{\frac\pi9,\frac{5\pi}9,\frac{7\pi}9\bigr\}$. Ces trois
    cosinus sont distincts ($\cos$ est injective sur $\intFF0\pi$) : ce sont
    les trois racines. Par les relations coefficients-racines, leur somme
    vaut $\res{0}$ (pas de terme en $t^{2}$) et leur produit $\res{\dfrac18}$.
    \pts{1,5 pt}
\end{questions}
\begin{center}\begin{minipage}{\linewidth}\centering
\begin{tikzpicture}[scale=2,line cap=round]
  \draw[otMuted,line width=0.5pt,-{Stealth[length=1.5mm]}] (-1.3,0) -- (1.35,0);
  \draw[otMuted,line width=0.5pt,-{Stealth[length=1.5mm]}] (0,-0.2) -- (0,1.25);
  \draw[otRule,line width=0.7pt] (1,0) arc (0:180:1);
  \foreach \a in {20,100,140}{
    \draw[otNavy,line width=0.7pt] (0,0) -- (\a:1);
    \draw[otGreenDark,line width=0.6pt,dash pattern=on 1.5pt off 1.2pt] (\a:1) -- ({cos(\a)},0);
    \fill[otNavy] (\a:1) circle (1.6pt);
    \fill[otGreenDark] ({cos(\a)},0) circle (1.6pt);}
  \node[font=\footnotesize,text=otNavy,anchor=south west] at (20:1.02) {$\pi/9$};
  \node[font=\footnotesize,text=otNavy,anchor=south] at (100:1.03) {$5\pi/9$};
  \node[font=\footnotesize,text=otNavy,anchor=south east] at (140:1.02) {$7\pi/9$};
  \node[font=\footnotesize,anchor=north] at (1,-0.03) {$1$};
  \node[font=\footnotesize,anchor=north] at (-1,-0.03) {$-1$};
\end{tikzpicture}\\*[0.5mm]
\legende{Les trois racines de $8t^{3}-6t-1$ sont les abscisses (en vert) des
points d'angles $\frac\pi9$, $\frac{5\pi}9$ et $\frac{7\pi}9$.}
\end{minipage}\end{center}
\end{corrigeEval}

\begin{corrigeEval}{6 points}{Nombres complexes}
\begin{questions}
  \item $\Delta=(3+\mathrm{i})^{2}-4(4+3\mathrm{i})=-8-6\mathrm{i}=(1-3\mathrm{i})^{2}$,
    d'où les racines
    $\dfrac{(3+\mathrm{i})\pm(1-3\mathrm{i})}{2}$ :
    $\res{2-\mathrm{i}\ \text{ et }\ 1+2\mathrm{i}}$. \pts{1 pt}
  \item On multiplie par le conjugué du dénominateur :
    \[Z=\frac{(z-\mathrm{i})\,\overline{(z+1)}}{|z+1|^{2}},\qquad
      (z-\mathrm{i})\,\overline{(z+1)}=\bigl(x^{2}+x+y^{2}-y\bigr)+\mathrm{i}\,(y-x-1).\]
    $Z$ réel $\iff y=x+1$ : c'est la \res{\text{droite passant par }-1\text{ et }\mathrm{i}},
    privée de $-1$. $Z$ imaginaire pur
    $\iff\bigl(x+\frac12\bigr)^{2}+\bigl(y-\frac12\bigr)^{2}=\frac12$ : c'est le
    \res{\text{cercle de diamètre }[-1,\mathrm{i}]}, privé de $-1$.
    \pts{2 pts}
  \item Pour $z=\mathrm{e}^{\mathrm{i}\theta}$, $\theta\notin2\pi\Z$, on
    factorise par l'angle moitié :
    \[\frac{1+z}{1-z}=\frac{\mathrm{e}^{\mathrm{i}\theta/2}\bigl(\mathrm{e}^{-\mathrm{i}\theta/2}+\mathrm{e}^{\mathrm{i}\theta/2}\bigr)}
      {\mathrm{e}^{\mathrm{i}\theta/2}\bigl(\mathrm{e}^{-\mathrm{i}\theta/2}-\mathrm{e}^{\mathrm{i}\theta/2}\bigr)}
      =\frac{2\cos\frac\theta2}{-2\mathrm{i}\sin\frac\theta2}
      =\res{\mathrm{i}\cot\frac\theta2}.\]
    \pts{1,5 pt}
  \item $z=1$ n'est pas solution ; on pose $w=\dfrac{z+1}{z-1}$, d'où
    $w^{n}=1$ et $w\neq1$ : $w=\mathrm{e}^{2\mathrm{i}k\pi/n}$ avec
    $1\leqslant k\leqslant n-1$, et $z=\dfrac{w+1}{w-1}$. D'après la
    question 3, $\res{z_k=-\mathrm{i}\cot\dfrac{k\pi}{n}}$ : $n-1$ solutions
    distinctes ($\cot$ est injective sur $\intOO0\pi$), toutes imaginaires
    pures. \pts{1,5 pt}
\end{questions}
\begin{center}\begin{minipage}{\linewidth}\centering
\begin{tikzpicture}[scale=1.25,line cap=round]
  \draw[otMuted,line width=0.5pt,-{Stealth[length=1.5mm]}] (-2,0) -- (1.2,0);
  \draw[otMuted,line width=0.5pt,-{Stealth[length=1.5mm]}] (0,-0.6) -- (0,1.7);
  \draw[otNavy,line width=1pt] (-1.6,-0.6) -- (0.6,1.6);
  \draw[otGreenDark,line width=1pt] (-0.5,0.5) circle (0.7071);
  \fill[otInk] (0,1) circle (1.5pt);
  \draw[otInk,line width=0.6pt,fill=white] (-1,0) circle (1.7pt);
  \node[font=\scriptsize,anchor=north east] at (-1,0) {$-1$};
  \node[font=\scriptsize,anchor=west] at (0.05,1.05) {$\mathrm{i}$};
  \node[font=\scriptsize,text=otNavy,anchor=west] at (0.62,1.6) {$Z\in\R$};
  \node[font=\scriptsize,text=otGreenDark,anchor=north] at (-0.5,-0.22) {$Z\in\mathrm{i}\R$};
\end{tikzpicture}\\*[0.5mm]
\legende{La droite et le cercle de la question 2, tous deux privés du point
$-1$ (point creux).}
\end{minipage}\end{center}
\end{corrigeEval}
